\documentclass[a4paper,11pt]{article}
\usepackage{exam-style-842-unit}
\usepackage{needspace}
\examinehead{842 信息化概论 · 第06单元练习册}{练习三\quad 物联网应用与新一代信息技术}
\begin{document}

\unitcover{842 信息化概论 · 单元练习册}{第06单元\quad 物联网与传感技术}{练习三\quad 物联网应用与新一代信息技术（共 100 分）}{约 70 分钟}{物联网典型应用（牲畜溯源·精准农业·安全防入侵·医疗保健）；停车场物联网案例四部分；智慧农业与农产品溯源；智能家居与智能交通；新一代信息技术 6 方面及其与物联网、云计算、大数据、人工智能的关系；安全与隐私；标准体系；产业前景。}

\begin{knowmap}
\textbf{应用图谱：}应用层面向绿色农业、工业监控、公共安全、城市管理、远程医疗、智能家居、智能交通、环境监测等行业建立管理与运行平台\\[2pt]
\textbf{典型应用：}牲畜溯源（RFID 电子耳码可查成长历史）· 精准农业（2002 年英特尔建立世界第一个无线葡萄园）· 安全防入侵（上海浦东国际机场铺设 3 万多个传感节点）· 医疗保健（可穿戴传感器实时回传健康参数）\\[2pt]
\textbf{停车场案例四部分：}感知设备（地磁传感器、摄像机＋车牌识别器、超声传感器）$\rightarrow$ 传输网络（3G/4G、TCP/IP）$\rightarrow$ 中心管理平台（通信、管理、计算、数据、媒体、应用服务器）$\rightarrow$ 监控中心（测控计算机＋电视大屏）\\[2pt]
\textbf{新一代信息技术 6 方面：}下一代通信网络（NGN）· 物联网 · 三网融合 · 新型平板显示 · 高性能集成电路 · 云计算\\[2pt]
\textbf{现状与挑战：}产业体系初步形成但产业化能力不高、核心关键技术待突破、标准比较分散、应用规模较小、安全与隐私保护挑战；大规模应用面临成本、安全、侵犯隐私 3 大挑战\\[2pt]
\textbf{标准化与前景：}国际 ETSI（M2M）· ITU-T（USN）· ISO/IEC（RFID）；我国 2005 年 RFID 标准工作组、2009 年传感器网络标准工作组（WGSN）、2010 年 CCSA 泛在网技术工作委员会、2010 年中国物联网标准联合工作组；EPoSS 预测物联网经历 4 个阶段，2020 年后全智能化，物物互联业务与人际通信业务之比将达 30:1
\end{knowmap}

\begin{keyterms}
\textbf{智慧农业}：以物联网、云计算、大数据、人工智能等新一代信息技术为支撑，实现农业生产全过程感知、决策、控制与管理的闭环，提高资源利用率和生产效益。\\
\textbf{精准农业}：基于农田、园林、温室等目标区域的射频技术与无线传感器网络架构，实时收集温度、湿度、光照、气体浓度等物理量，精准获取土壤水分、压实程度、电导率、pH 值、氮素等信息，实现精细化的农业生产控制。\\
\textbf{农产品溯源}：利用 RFID、二维码等标识技术记录农产品种养、加工、流通全过程信息，实现来源可查、去向可追、责任可究。\\
\textbf{电子耳码}：给牲畜佩戴的 RFID 电子标签，记录牲畜身份与成长历史，可一直保留到超市出售的肉品上供消费者查询。\\
\textbf{无线葡萄园}：2002 年英特尔在俄勒冈建立的世界上第一个无线葡萄园，节点每分钟检测土壤温度、湿度或有害物数量。\\
\textbf{停车场物联网}：由感知设备、传输网络、中心管理平台和监控中心 4 部分组成的典型物联网应用。\\
\textbf{中心管理平台}：部署在互联网云计算平台或停车场业主机房内，由通信服务器、管理服务器、计算服务器、数据服务器、媒体服务器和应用服务器等组成。\\
\textbf{新一代信息技术}：包括下一代通信网络、物联网、三网融合、新型平板显示、高性能集成电路和云计算 6 个方面。\\
\textbf{下一代通信网络}：建立在高带宽光纤网络和 3G 无线网络基础上，实现音频、视频、数据统一的综合通信网络。\\
\textbf{三网融合}：电信网、移动互联网以及广播电视网的融合，是应用上的有机融合而非物理网络的合一。\\
\textbf{新型平板显示}：以 OLED 和电子纸为代表，具有显示效果好、轻薄省电、可柔性弯折或可卷折等特点。\\
\textbf{云计算}：将计算任务分布在由大规模数据中心或大量计算机集群构成的资源池上，使用户按需获取计算能力、存储空间和软件服务的技术。\\
\textbf{物联网标准}：保证物联网产品互操作性和互联互通的规范体系，目前国际国内均未形成统一标准，是产业发展的关键。\\
\textbf{隐私保护}：物联网承载大量经济社会活动和战略性资源，采集个人与组织数据时必须防止信息被滥用、泄露和侵犯隐私。
\end{keyterms}

\needspace{6cm}
\begin{contrast}{智慧农业 / 精准农业（每行依次为：核心内涵·关键特征·支撑技术）}
\citem{智慧农业}{以新一代信息技术支撑农业生产全过程感知、决策、控制与管理}{强调“感知—决策—控制”闭环，覆盖生产、加工、流通、销售全链条}{物联网、云计算、大数据、人工智能、3S 技术}
\citem{精准农业}{基于目标区域的射频技术与无线传感器网络架构，精细化控制农业生产}{强调按需、定位、定量投入，实时获取土壤与气象信息}{传感器、无线传感器网络、GPS、GIS、遥感}
\end{contrast}

\needspace{6cm}
\begin{contrast}{物联网 / 云计算 / 大数据 / 人工智能（每行依次为：技术定位·核心功能·在农业物联网中的作用）}
\citem{物联网}{新一代信息技术的感知与连接基础}{通过传感器、RFID 等全面感知并传输物物信息}{负责“采得到、传得回”，是数据源头}
\citem{云计算}{新一代信息技术的计算与存储基础}{把计算任务分布到资源池，按需提供计算、存储和软件服务}{负责“存得下、算得动”，是支撑层核心}
\citem{大数据}{新一代信息技术的数据处理基础}{对海量、多源、异构数据进行存储、分析与挖掘}{负责“挖得深”，从数据中发现规律}
\citem{人工智能}{新一代信息技术的智能决策基础}{用计算机模拟、延伸人的智能行为和思维过程}{负责“判得准”，实现病虫害识别、灌溉决策等}
\end{contrast}

\needspace{6cm}
\begin{contrast}{停车场物联网四部分（每行依次为：主要设备·承担功能·数据流向）}
\citem{感知设备}{地磁传感器、摄像机、车牌识别器、超声传感器、信号灯}{感知车辆出入与车位占用情况}{采集车辆、车位、时间等信息}
\citem{传输网络}{3G/4G 移动通信网、互联网、TCP/IP}{把分布在各处的感知设备接入互联网平台}{感知数据上传至中心管理平台}
\citem{管理平台}{通信、管理、计算、数据、媒体、应用服务器}{计算处理、存储数据、提供业务服务}{向下接收数据，向上支撑监控中心}
\citem{监控中心}{测控计算机、电视大屏}{为停车场运营管理和安全保障提供决策信息}{展示平台处理结果，供管理者调度}
\end{contrast}

\qtypename{一、单项选择题}{10 题，每题 3 分，共 30 分}
\q 在牲畜溯源中，给牲畜佩戴的 RFID 电子标签俗称（\quad）
\choicesline{A. 电子耳码}{B. 二维码}{C. 电子鼻}{D. 电子秤}

\q 2002 年率先建立世界上第一个无线葡萄园的公司是（\quad）
\choicesline{A. 英特尔}{B. IBM}{C. 微软}{D. 谷歌}

\q 上海浦东国际机场防入侵系统铺设的传感节点数量为（\quad）
\choicesline{A. 3 千多个}{B. 3 万多个}{C. 30 多万个}{D. 300 多个}

\q 停车场案例中，出入口地磁传感器感知到车辆出入行为后立即启动的设备是（\quad）
\choices{A. 道闸}{B. 摄像机与车牌识别器}{C. 信号灯}{D. 电视大屏}

\q 停车场案例中各车位超声传感器采用的短距离无线通信组网技术是（\quad）
\choicesline{A. ZigBee}{B. 卫星通信}{C. 光纤通信}{D. 广播电视网}

\q 三网融合主要指（\quad）
\choices{A. 三张物理网络合并成一张}{B. 电信网、移动互联网、广播电视网的应用融合}
{C. 三家企业合并经营}{D. 三种终端的融合}

\q 下列不属于新一代信息技术的是（\quad）
\choicesline{A. 物联网}{B. 三网融合}{C. 云计算}{D. 传统机械制造}

\q 下列不属于教材所述物联网大规模应用面临的挑战的是（\quad）
\choicesline{A. 成本的挑战}{B. 安全的挑战}{C. 侵犯隐私的威胁}{D. 标准完全统一}

\q 教材指出，家庭自动化（智能家居）要给人创造的居住环境是（\quad）
\choices{A. 安全、高效、舒适、节能、环保}{B. 豪华、昂贵、气派}
{C. 封闭、独立、自给}{D. 单一、固定、低耗}

\q 预测物联网发展将经历 4 个阶段的组织及其报告是（\quad）
\choices{A. 欧盟 EPoSS《Internet of Things in 2020》}{B. 美国 FORRESTER 公司报告}
{C. ITU《ITU 互联网报告 2005：物联网》}{D. 中投顾问《中国物联网行业调研报告》}

\qtypename{二、多项选择题}{4 题，每题 4 分，共 16 分；每题有两个或两个以上正确项}
\q 教材 7.1.3 所列的物联网典型应用包括（\quad）
\choicesline{A. 牲畜溯源}{B. 精准农业}{C. 安全防入侵}{D. 数控机床加工}

\q 我国物联网发展现状存在的问题有（\quad）
\choices{A. 产业化能力不高}{B. 核心关键技术待突破}
{C. 标准比较分散、体系不完善}{D. 应用规模大、商业模式成熟}

\q 下列属于新一代信息技术的有（\quad）
\choicesline{A. 下一代通信网络}{B. 物联网}{C. 三网融合}{D. 传统机械制造}

\q 物联网产业发展面临的挑战包括（\quad）
\choicesline{A. 标准完全统一}{B. 成本的挑战}{C. 安全的挑战}{D. 侵犯隐私的威胁}

\qtypename{三、判断题}{5 题，每题 2 分，共 10 分；正确的画“√”，错误的画“×”并写出改正后的表述}
\q 精准农业是指研究基于农田、园林、温室等目标区域的射频技术与无线传感器网络架构，实现实时收集温度、湿度、光照、气体浓度等物理量。\quad（\quad）\par\noindent 改正：\fixblank

\q 物联网的应用层是物联网系统结构的最高层，包括各类用户界面显示设备及其他管理设备。\quad（\quad）\par\noindent 改正：\fixblank

\q 我国物联网标准体系已经完善，国际上也已经形成统一的物联网标准。\quad（\quad）\par\noindent 改正：\fixblank

\q 三网融合就是电信网、移动互联网以及广播电视网在物理网络上的合一。\quad（\quad）\par\noindent 改正：\fixblank

\q 智慧农业就是在农田里安装摄像头进行远程监控。\quad（\quad）\par\noindent 改正：\fixblank

\qtypename{四、配对题}{4 对，每对 1 分，共 4 分；把 A 栏与 B 栏对应内容配对}
\q 将下列应用与其对应内容配对，把 B 栏字母填入括号。
\par\vspace{0.2em}
\noindent\begin{tabular}{@{}p{0.44\textwidth}@{\hspace{1.2em}}p{0.44\textwidth}@{}}
（1）精准农业 \hfill（\quad） & A. 电子耳码 \\
（2）牲畜溯源 \hfill（\quad） & B. 无线葡萄园 \\
（3）安全防入侵 \hfill（\quad） & C. 通信、管理、计算、数据服务器等 \\
（4）停车场中心管理平台 \hfill（\quad） & D. 上海浦东国际机场 3 万多个传感节点 \\
\end{tabular}\par\vspace{0.3em}

\newpage
\qtypename{五、简答题}{4 题，每题 6 分，共 24 分}
\q 简述物联网在农业中的两个典型应用（牲畜溯源、精准农业）及其技术实现。\anslines{2.4cm}

\q 简述停车场物联网案例的四个组成部分及各自作用。\anslines{2.4cm}

\q 简述新一代信息技术的 6 个方面，并说明物联网与云计算、大数据、人工智能之间的关系。\anslines{2.4cm}

\q 简述我国物联网发展的现状、标准化概况、产业前景及面临的主要挑战。\anslines{2.4cm}

\newpage
\qtypename{六、综合题}{1 题，共 16 分}
\q 某县建设“智慧农业＋农产品溯源”平台：稻田部署小型气象站和土壤传感器，农机安装北斗终端，稻米加工包装环节赋二维码和 RFID 标签；消费者扫码可查产地、农事记录和检测报告；平台依托云计算存储数据，用大数据分析产量影响因素，用人工智能模型识别病虫害并推送防治方案；县里同时出台地方标准、建设数据管理制度。请回答：（1）说明物联网在该县稻米“种植—加工—物流—销售”各环节中的具体作用；（2）分析云计算、大数据、人工智能在平台中各自承担的角色及其协同关系；（3）从安全与隐私保护、标准体系两个角度，提出 3 条改进建议。\anslines{5cm}

\end{document}
